% =============================================================================
% Appendix A: Ritz Basis and Convergence
% =============================================================================
% PURPOSE:
%   Make Paper 10 self-contained by stating the oblate spheroidal Ritz model
%   explicitly: coordinate system, trial functions, energy assembly, boundary
%   conditions, n=1 treatment, and convergence evidence.
% =============================================================================

\appendix
\section{Oblate Spheroidal Shell Ritz Model}
\label{sec:appendix_ritz}

This appendix collects the Rayleigh--Ritz formulation used throughout the
companion papers and the present capstone.  The model follows the classical
shell-acoustics tradition~\cite{lamb1882,Junger1986} with a Love--Kirchhoff
thin-shell assumption.

\subsection{Coordinate system and trial functions}
\label{sec:appendix_coordinates}

The undeformed shell midsurface is an oblate spheroid with semi-major axis $a$
and semi-minor axis $c$ ($c < a$).  Define the eccentricity
$\varepsilon = \sqrt{1-c^2/a^2}$ and the semifocal distance
$d = a\varepsilon$.  In oblate spheroidal coordinates $(\xi,\mu,\varphi)$, the
shell surface sits at $\xi_0 = c/d$.  Here $\mu = \cos\theta$ ranges over
$[-1,1]$ with $\theta$ the polar angle, and $\varphi\in[0,2\pi)$ is the
azimuthal angle.  (We write $\mu$ rather than the traditional $\eta$ for this
coordinate in order to avoid conflict with the material loss tangent~$\eta$
used elsewhere in the paper.)  The metric function
\begin{equation}
  G(\mu) = a^2\mu^2 + c^2(1-\mu^2)
  \label{eq:app_metric}
\end{equation}
determines the position-dependent curvature.  The surface element is
$\mathrm{d}S = 2\pi\,a\sqrt{G}\,\mathrm{d}\mu$.

For each modal order $n\geq 2$, the trial displacement field
is a two-degree-of-freedom (2-DOF) Ritz ansatz:
\begin{equation}
  w(\mu) = \beta\,P_n(\mu), \qquad
  u(\theta) = \alpha\,\sin\theta\,P_n'(\cos\theta),
  \label{eq:app_trial}
\end{equation}
where $w$ is the normal (radial) displacement, $u$ is the tangential
(meridional) displacement, $P_n$ denotes the Legendre polynomial of degree~$n$,
and $\beta,\alpha$ are the Ritz amplitudes forming the generalised coordinate
vector $\mathbf{q}=(\beta,\alpha)^\top$.  Legendre polynomials are smooth
on $[-1,1]$ and regular at both poles, so the trial functions are admissible
for a free (unconstrained) shell without requiring any additional boundary
conditions.

\subsection{Energy functionals and eigenvalue problem}
\label{sec:appendix_energy}

The stiffness matrix $\mathbf{K}$ receives three contributions, all integrated
over the shell surface via Gauss--Legendre quadrature (200 nodes by default):
\begin{enumerate}
  \item \emph{Membrane stiffness.}  With extensional rigidity
    $C_m = Eh/(1-\nu^2)$, the membrane strain energy arises from meridional
    and circumferential strains that depend on $P_n$, $P_n'$, and the
    spatially varying metric~\eqref{eq:app_metric}.
  \item \emph{Bending stiffness.}  With flexural rigidity
    $D_f = Eh^3/[12(1-\nu^2)]$, the curvature-change energy involves $P_n''$
    and the derivatives of the metric function, weighted by $\sqrt{G}$.
  \item \emph{Geometric prestress.}  The intra-abdominal pressure
    $P_\mathrm{iap}$ induces a membrane stress resultant
    $N_1 = P_\mathrm{iap}\,R_2/2$, where $R_2(\mu) = a\sqrt{G}/c$ is the
    circumferential radius of curvature.  This contributes a geometric
    stiffness term (tension stiffening) that acts on the normal DOF only.
\end{enumerate}

The mass matrix $\mathbf{M}$ has two contributions:
\begin{enumerate}
  \item \emph{Shell inertia.}  The structural mass per unit area is
    $\rho_w h$, integrated against $P_n^2$ for the normal DOF and
    $\sin^2\theta\,(P_n')^2$ for the tangential DOF.
  \item \emph{Fluid added mass.}  The interior fluid loading is computed from
    the Laplace problem in oblate spheroidal coordinates using Legendre
    functions $P_n(i\xi)$ on the imaginary argument, following Junger and
    Feit~\cite{Junger1986}.  This yields a scalar added-mass contribution to
    the normal DOF.  Near the spherical limit ($c/a > 0.98$), the oblate
    harmonic solution is blended with the Lamb sphere formula via a
    $C^2$-continuous interpolant to avoid numerical singularity in the
    coordinate transformation.
\end{enumerate}

Stationarity of the Rayleigh quotient yields the $2\times 2$ generalised
eigenvalue problem
\begin{equation}
  \mathbf{K}\mathbf{q} = \omega^2\mathbf{M}\mathbf{q},
  \label{eq:app_eigenproblem}
\end{equation}
from which the natural frequency of mode $n$ is extracted as the lowest
positive eigenvalue.  Solving \eqref{eq:app_eigenproblem} independently for
each $n$ produces the flexural spectrum $\{f_n\}_{n\geq 2}$ used throughout
this paper.

\subsection{Free boundary conditions and rigid-body modes}
\label{sec:appendix_bcs}

The shell is modelled as free (unconstrained).  No displacement or rotation
boundary conditions are imposed.  For $n=1$, the deformation is a rigid-body
translation of the shell; the fluid adds mass but provides no restoring force,
so $\omega_1=0$.  The present paper considers only $n\geq 2$, for which the
restoring force arises from shell elasticity and prestress.

\subsection{Convergence of the Ritz basis}
\label{sec:appendix_convergence}

The 2-DOF trial~\eqref{eq:app_trial} can be enriched by retaining additional
Legendre terms of the same parity:
\begin{equation}
  w(\mu) = \sum_{i=0}^{N-1}\beta_i\,P_{n+2i}(\mu), \qquad
  u(\theta) = \sum_{j=0}^{N-1}\alpha_j\,\sin\theta\,P_{n+2j}'(\cos\theta),
  \label{eq:app_enriched}
\end{equation}
giving a $2N$-DOF problem.  Note that in the enriched basis, the fluid
added mass is computed per-degree using oblate spheroidal harmonics for each
individual Legendre order $n+2i$; off-diagonal fluid coupling between
different Legendre orders is neglected.

For the canonical abdominal parameters (Table~1) and the $n=2$ mode, the
convergence study in~\cite{BrowntoneP1} yields the results summarised in
Table~\ref{tab:app_convergence}.

\begin{table}[h]
  \centering
  \small
  \caption{Convergence of the $n=2$ flexural frequency with basis enrichment.
    $N$ is the number of Legendre terms per displacement component;
    $\Delta f$ is the relative deviation from the $N=8$ reference value.}
  \label{tab:app_convergence}
  \renewcommand{\arraystretch}{1.1}
  \begin{tabular}{@{}ccc@{}}
    \toprule
    $N$ (DOFs = $2N$) & $f_2$ (Hz) & $\Delta f$ vs.\ $N{=}8$ \\
    \midrule
    1 \phantom{0}(2)  & 3.800 & $+3.3\%$ \\
    2 \phantom{0}(4)  & 3.682 & $+0.1\%$ \\
    3 \phantom{0}(6)  & 3.679 & $<0.01\%$ \\
    $\geq 4$ \phantom{0}($\geq 8$) & 3.678 & converged \\
    \bottomrule
  \end{tabular}
\end{table}

Three properties confirm adequacy of the 2-DOF basis.  First, frequencies
decrease monotonically with $N$, consistent with the Rayleigh--Ritz
upper-bound property.  Second, the 2-DOF model ($N=1$) overestimates the
converged frequency by only \SI{3.3}{\percent}, and the 4-DOF model ($N=2$)
is within \SI{0.1}{\percent}.  Third, the canonical five-mode set
($n=2,\ldots,6$) used for the identifiability analysis lies well above the
convergence threshold.

To verify that the inverse-problem quantities are stable with respect to the
number of observed modes, we computed the condition number $\kappa$ for
mode sets of increasing size:
$\kappa = 73.0$ (3 modes), $66.7$ (4), $69.4$ (5), $74.4$ (6), $80.5$ (7).
The third singular value $\sigma_3$ likewise stabilises:
$\sigma_3 = 0.019$ (3 modes), $0.024$ (4), $0.026$ (5), $0.027$ (6), $0.028$ (7).
For five or more modes, $\sigma_3$ varies by less than $6\%$.

\medskip\noindent
\textbf{Basis enrichment.}
Table~\ref{tab:app_convergence} demonstrates that the single-term Ritz basis
overestimates the converged $n=2$ frequency by only \SI{3.3}{\percent}.
The structural properties of the identifiability analysis --- rank collapse at
the sphere, even symmetry in $\varepsilon$, and the consequent
$\sigma_3\sim\varepsilon^2$ scaling --- are algebraic features of the model
family that hold irrespective of basis order $N$, since they follow from
symmetry and spectral perturbation theory rather than from particular
numerical values.  The specific numerical value $\kappa=69.4$ may shift
modestly with basis enrichment, but the four-orders-of-magnitude gap between
$\kappa_\mathrm{oblate}\approx 70$ and
$\kappa_\mathrm{sphere}\approx 10^{10}$ ensures that the identifiability
lifting conclusion is robust against the \SI{3.3}{\percent} frequency bias
inherent in the single-term basis.
